\label{sec:gates}

This section specifies how a candidate marker earns a reliability tier -- the machinery applied
uniformly to all seven markers (six prespecified, one post hoc; Section~\ref{sec:prespecified}).

\subsection{Fixed evaluation protocol}

All continuous probes use \texttt{StandardScaler + RidgeCV(alphas=logspace(-2,6,25))}; all
binary probes use \texttt{StandardScaler + LogisticRegression(C=1.0, class\_weight="balanced",
max\_iter=2000)} -- one fixed protocol across every marker and cohort, adopted after finding
that per-test hyperparameter search could flip a borderline result (marker 3's TCGA-PRAD
ERG-fusion-status replication moved from nominally significant to null once standardized;
Section~\ref{sec:failure-modes}).

\subsection{Cross-validation and statistical protocol}

All within-cohort evaluation uses patient-disjoint \texttt{GroupKFold(n\_splits=5)} (a
patient's slides never split across train/test), motivated by an early finding that our own
rule-based algorithm's apparent slide-level signal on NADT-Prostate was pseudo-replication that
disappeared at the patient level. Primary metrics are patient-level (mean prediction vs.\ mean
or majority true label per patient); slide-level metrics are reported as secondary. Across the
13 distinct marker-hypothesis tests conducted on real ground truth in this project (the 6
initial frozen-pool markers plus 7 additional TCGA-PRAD molecular candidates that were not
promoted to the pool) and the four nested/refit incremental-value tests introduced by the
confounder audit, Benjamini--Hochberg FDR correction is applied as one 17-test family;
external-cohort re-tests and cross-encoder replications of an already-tested hypothesis are
excluded from this family, since they confirm rather than newly test a hypothesis. Confidence
intervals are patient-ID-cluster bootstrap (2{,}000 resamples).

\subsection{Confounder audit (Gate 4)}
\label{sec:gate4-methods}

For markers whose target is plausibly confounded by grade (4, 6, and -- retrospectively -- 7),
three models are compared: clinical-only, image-only, and combined. For markers 4/6, each
outer patient-disjoint fold constructs the meta-model's training image feature from inner-fold
OOF predictions, refits the image probe on the complete outer-training fold, and applies both
probe and meta-model to untouched outer-test patients. Marker 7's image score is fixed because
it was trained externally on LEOPARD; only clinical and combined Cox coefficients are fit in
each TCGA-PRAD outer fold. Primary statistics are held-out patient-level $\Delta$AUROC,
$\Delta R^2$, or $\Delta$C-index with 2{,}000 patient-bootstrap confidence intervals. The
earlier full-cohort likelihood-ratio analyses are retained as auxiliary, explicitly
in-sample results (Section~\ref{sec:prespecified} discusses why this stricter nested design
was adopted after the original results).

The nonparametric check permutes patient labels within rounded patient-mean Gleason strata and
reruns the complete nested pipeline, including image-probe refitting for markers 4/6, for
2{,}000 repetitions. For marker 7, paired event/time outcomes are permuted within grade
strata; its externally trained image score remains fixed while all outer-fold Cox models are
refit. Markers 1 and 3 are excluded from this audit because their target is grade itself,
making a ``does the image add information beyond grade'' test circular.

\subsection{Reliability tiers}

Each initial marker is assigned one of five prospectively frozen tiers, from strongest to
weakest evidence: \emph{externally transportable} (a fixed, source-cohort-fitted probe applied
with no retraining reproduces the association in an independent cohort); \emph{cross-cohort
replicable} (the association reproduces after refitting in an independent cohort or a
within-cohort held-out site, but no zero-shot transport test exists); \emph{internally
supported, externally untested} (internally consistent across encoders and/or molecular
assays, but no independent cohort with the required label exists, confirmed by an exhaustive
dataset search); \emph{context-sensitive} (direction or magnitude is unstable across site,
scale, encoder, or assay); \emph{unsupported} (does not meet the specified statistical support
criteria under the tested design). The historical protocol label was \emph{unsupported/null};
here, ``unsupported'' means statistical non-support under the tested design, not proof of a
biological null. Marker 7, discovered post hoc, is described using this same vocabulary in
Section~\ref{sec:recurrence-transfer} without being assigned a prospective tier.

\subsection{Gate A sensitivity-grid summaries}

The completed frozen grid contains 360 correlated seed-level cells: six markers, two encoders,
tile budgets 16/32/64, five tile-sampling seeds, and encoder-specific native versus shared
$1.76\,\mu$m/px scales. AUROC and C-index cells are chance-or-worse at $\leq0.5$; Spearman
$\rho$ cells at $\leq0$. The 72 configuration rows report the five-seed mean, sample SD,
range, chance-or-worse count, and a df=4 Student-$t$ interval across sampling seeds. These
intervals describe repeated tile sampling on the same frozen patients/folds, not patient-level
or population uncertainty. The 390 paired contrasts preserve each seed-level pairing and
report null crossings for native versus $1.76\,\mu$m/px, Virchow versus CONCH at shared scale,
and 64 versus 16 tiles; no independence-based $p$-values are computed. Because settings share
cohorts and folds, this is a descriptive correlated sensitivity audit rather than 360
confirmatory tests, external validation, or independent replication. Reliability tiers are
not elevated by the grid alone. The Gate B-approved framing treats the grid as a bound on
generalizability, not as evidence that prior directional observations were overturned; open
questions remain unresolved rather than contrary evidence. R5 evidence closure for C04/C05 is
complete, while their scientific questions remain unresolved. C06 has completed the
common-498 R2 and official-PFI R3 endpoint sensitivities. C07 is closed narrowly by the
same-patient/same-draw R4 analysis: the endpoint-conditioned results do not establish robust
independent or incremental prognostic value.

\subsection{QWK conversion}

To compare marker 1 against literature figures reported as F1 or quadratic weighted kappa
(QWK) rather than Spearman correlation, continuous predictions are converted to ordinal
categories via quantile binning matched to the number of true grade categories present in each
cohort (a transparent conversion that does not use true-label thresholds beyond the category
count), and QWK is computed against the true ordinal grade. This is applied uniformly across
all five evaluated cohorts.
